\Def Сортировка набора элементов $a_1, \dots , a_n$ — задача поиска такой
перестановки индексов $\sigma$, что {\large $a_{\sigma(1)} \leq a_{\sigma(2)} \leq \, \dots \, \leq a_{\sigma(n)}$}.

\Lemma $\log{n!}=\Theta(n\log{n})$\\
\Proof
$$\log{n!}=\log{1}+\log{2}+\ldots+\log{n} \leq n\log{n} \Rightarrow \log{n!}=O(n\log{n})$$
$$\log{n!} \geq \frac{n}{2}\cdot \log{\frac{n}{2}} = \frac{n}{2}(\log{n} - \log{2}) = \frac{1}{2}n\log{n} - \frac{1}{2}n\log{2} \Rightarrow \log{n!}=\Omega(n\log{n})$$
Следовательно, $\log{n!}=\Theta(n\log{n})$ \Endproof

\Th Сортировка N элементов, основанная на сравнениях, требует $\Omega(N\,log\,N)$ сравнений.

\Proof \\ Задача состоит в поиске одной из $n!$ перестановок. Заметим, что при ответе на запрос: какое из двух чисел должно стоять раньше? - количество перестановок уменьшается в два раза. То есть $n! \mapsto \frac{n!}{2} \mapsto \frac{n!}{4} \mapsto \ldots \;$, что влечет хотя бы $\log{n!}$ итераций $\Rightarrow$ (по лемме) $\Omega(n\log{n})$. \Endproof

\Def Сортировка стабильная, если одинаковые с точки зрения компаратора элементы не меняют своего относительного расположения.

Из рассмотренных сортировок ранее, только MergeSort стабильная

\pagebreak